\documentclass[12pt]{article}

\usepackage[utf8]{inputenc}
\usepackage{amsmath}
\usepackage{amssymb}
\usepackage{graphicx}
\usepackage{bbm}
\usepackage{jheppub}

\newcommand{\Ident}{\mathbbm{1}}
\newcommand{\svx}{\mathsf{x}}

\title{XXX CG}

\author{Paul Ryan}

\emailAdd{paul.ryan@desy.de}

\abstract{This is the start of the xxx-CG document.}

\begin{document}

\maketitle

\section{Introduction}

This is the start of the xxx-CG document.

\section{The XXX spin chain and separation of variables}
\label{sec:xxx-sov}

Before addressing the Gaudin model, we first set up the closely related rational $\mathfrak{gl}(2)$ (XXX) spin chain, for which separation of variables (SoV) is completely explicit. The Gaudin model statements we are ultimately after are recovered from the XXX spin chain in an appropriate scaling limit.

\subsection{Lax matrix and transfer matrix}

Let $V_\lambda$ denote the finite-dimensional irreducible representation of $\mathfrak{gl}(2)$ with highest weight $(\lambda,0)$, of dimension $\lambda+1$, realized on the Gelfand--Tsetlin basis. Denote by $E_{ij}\in\mathrm{End}(V_\lambda)$, $i,j=1,2$, the representation matrices of the $\mathfrak{gl}(2)$ generators.

The Lax matrix in the fundamental (two-dimensional) auxiliary space and physical space $V_\lambda$ is
\[
L^{[\lambda]}(u) = u\,\Ident - h\sum_{i,j=1}^2 E_{ji}\otimes e_{ij} \;\in\; \mathrm{End}(V_\lambda)\otimes\mathrm{End}(\mathbb{C}^2),
\]
where $e_{ij}$ are the elementary matrices acting on the auxiliary space and $h$ is the (quasi-classical) parameter that will later be scaled to zero in passing to the Gaudin model.

For a single inhomogeneous site ($L=1$) with inhomogeneity $\theta$ and twist $G\in GL(2)$, the transfer matrix is
\[
t_{1,1}(u) = \mathrm{tr}_a\big[G\,L^{[\lambda]}(u-\theta)\big] .
\]

\subsection{Companion twist and the fusion hierarchy}

It is convenient to parametrize the twist $G$, with eigenvalues $\kappa_1,\kappa_2$, in companion form
\[
G = \begin{pmatrix} \chi_1 & -\chi_2 \\ 1 & 0 \end{pmatrix}, \qquad
\chi_1=\kappa_1+\kappa_2,\quad \chi_2=\kappa_1\kappa_2 ,
\]
rather than diagonal form. This choice makes the SoV construction below algebraically trivial.

Higher transfer matrices, associated with auxiliary representations labeled by a rectangular Young diagram with $a$ rows and $s$ columns, are denoted $t_{a,s}(u)$. The antisymmetric fusion $a=2,\,s=1$ reduces to the (twisted) quantum determinant,
\[
t_{2,1}(u) = \chi_2\, \mathrm{qdet}\, L^{[\lambda]}(u-\theta), \qquad
\mathrm{qdet}\,L^{[\lambda]}(u) = L^{[\lambda]}_{11}(u)L^{[\lambda]}_{22}(u-h) - L^{[\lambda]}_{21}(u)L^{[\lambda]}_{12}(u-h),
\]
and the symmetric fusions $t_{1,s}(u)$ are generated from $t_{1,1}$ and $t_{2,1}$ by the Hirota (T-system) recursion
\[
t_{1,s}(u) = t_{1,s-1}(u)\,t_{1,1}\big(u+(s-1)h\big) - t_{2,1}\big(u+(s-1)h\big)\,t_{1,s-2}(u),
\]
with the boundary conditions $t_{0,s}(u)\equiv\Ident$ and $t_{a,0}(u)\equiv\Ident$ for all $a,s\geq0$.

\subsection{Separation of variables}

Define $Q_\theta(u) = u-\theta$. In the companion-twist gauge the Sklyanin SoV $B$-operator is simply the $(1,1)$ matrix element of the Lax matrix,
\[
B^{[\lambda]}(u) = L^{[\lambda]}_{11}(u-\theta) ,
\]
so its spectrum is read off directly without any further diagonalization.

Let $\langle \svx^{[\lambda]}_0|\in V_\lambda^*$ denote the lowest-weight covector of the Gelfand--Tsetlin basis; this is the SoV pseudovacuum. Acting on it with the fused transfer matrices generates the full SoV covector basis,
\[
\langle \svx^{[\lambda]}_s| \;:=\; \frac{\langle \svx^{[\lambda]}_0|\, t_{1,s}(\theta)}{\displaystyle\prod_{p=1}^{s} Q_\theta\big(\theta-\lambda h+(p-1)h\big)}, \qquad s=0,1,\dots,\lambda ,
\]
and one checks that these are left eigenvectors of the $B$-operator,
\[
\langle \svx^{[\lambda]}_s|\, B^{[\lambda]}(u) = (u-\theta-sh)\, \langle \svx^{[\lambda]}_s| .
\]
This has been verified explicitly, symbolically in $u$, for all $\lambda=1,\dots,5$ and $s=0,\dots,\lambda$. Since $s$ ranges over $0,\dots,\lambda$, the $\lambda+1$ covectors $\langle \svx^{[\lambda]}_s|$ exhaust the SoV basis of $V_\lambda^*$ at $L=1$; for $\lambda=\tfrac12,1$ they were additionally checked to reproduce the Gelfand--Tsetlin basis exactly.

\bibliographystyle{JHEP}
\bibliography{references}

\end{document}
